\documentstyle[12pt,cours]{article}
%\largerfonts
\begin{document}

%%%%%%%%%%%%%%%%%%%%%%%%%%%%%%%%%%%%%%%%%%%%%%%%%%
%%%%%%%%%%%%%%%%%%

\begin{slide}
\begin{center}
{\Huge \bf Les Mod\`eles Logiques de Donn'ees\\}
\vspace*{1cm}
{\Large Philippe RIGAUX\\}
\vspace*{1cm}
{\large rigaux@cnam.fr \\}
\vspace*{1cm}

\end{center}

\end{slide}

%%%%%%%%%%%%%%%%%%%%%%%%%%%%%%%%%%%%%%%%%%%%%%%%%%

% Introduction 1

\begin{slide}
\subsection{Contenu}

\begin{itemize}
  \item Repr\'esentation physique~: fichiers et bases de donn\'ees
  \item Repr\'esentation conceptuelle, logique et physique
  \item Du MCD au mod\`ele hi\`erarchique
  \item Du MCD au mod\`ele r\'eseau
  \item Du MCD au mod\`ele relationnel
  \item Retour sur quelques probl\`emes~: choix des identifiants,
	optimisation, ...
\end{itemize}

\end{slide}

%%%%%%%%%%%%%%%%%%%%%%%%%%%%%%%%%%%%%%%%%%%%%%%%%%

\section{Repr\'esentation physique~: fichiers et bases de donn\'ees}

\begin{slide}[1.5cm]
\subsection{Qu'est-ce qu'une base de donn\'ees}

C'est un ensemble {\bf important} de donn\'ees avec les
caract\'eristiques suivantes~:

\begin{itemize}
  \item {\bf Permanence}~: les donn\'ees
        doivent survivre aux programmes,
  \item {\bf Structuration}~: on doit transcrire (et stocker) au niveau
	physique les notions conceptuelles
	d'entit\'es, de relations, de propri\'et\'es, de cardinalit\'es.
  \item {\bf Accessibilit\'e}~: on doit pouvoir extraire des
	informations et leur structure. 
\end{itemize}

$\rightarrow$ une base de donn\'ees est {\bf toujours} implant\'ee
		dans un ou plusieurs fichiers
\end{slide}

%%%%%%%%%%%%%%%%%%%%%%%%%%%%%%%%%%%%%%%%%%%%%%%%%%

% 2 Introduction 2

\begin{slide}[1.5cm]
\subsection{Rappels sur les fichiers}

\begin{itemize}
 \item Un fichier = un ensemble de blocs sur un disque
 \item Organisation en {\bf articles}
 \item Un article = 1 liste de champs
\end{itemize}

\noindent
Fonctions sur les fichiers~:

\begin{itemize}
  \item Insertion
  \item Suppression
  \item Modification
  \item Recherche
\end{itemize}

%%%%%%%%%%%%%%%%%%%%%%%%%%%%%%%%%%%%%%%%%%%%%%%%%%

\begin{slide}
\subsection{Une application bas\'ee sur un SGF}
%\centerline{\psfig{height=\textheight,figure=appli-sgf.ps}}
\end{slide}

%%%%%%%%%%%%%%%%%%%%%%%%%%%%%%%%%%%%%%%%%%%%%%%%%%

\begin{slide}[1.5cm]
\subsection{Inconv\'enients d'un SGF}

Chaque programme est responsable individuellement de
la pr\'eservation de la base. Or les probl\`emes suivants~:
peuvent survenir.

\begin{itemize}
  \item Les s\'equences d'acc\`es au donn\'ees sont dupliqu\'ees
	dans chaque fichier.
  \item Evolution d'un fichier = \'evolution de tous les programmes
  \item Concurence d'acc\`es~: deux programmes modifient en m\^eme
	temps les m\^emes donn\'ees.
  \item S\'ecurit\'e~: un programme d\'efectueux peut corrompre la
	base.
\end{itemize}
	
\noindent
Solution~: un syst\`eme interm\'ediaire, le {\bf Syst\`eme de Gestion
de Base de Donn\'ees (SGBD)} prend en charge l'acc\`es au fichier.
\end{slide}

%%%%%%%%%%%%%%%%%%%%%%%%%%%%%%%%%%%%%%%%%%%%%%%%%%

\begin{slide}[1.5cm]
\subsection{Les fonctions d'un SGBD}

Un SGBD propose les outils suivants~:

\begin{itemize}
  \item {\bf Mod\`eles de donn\'ees}~: description de l'organisation des
      donn\'ees, dit {\bf Mod\`ele logique}.

  \item {\bf S\'ecurit\'e}~: contr\^ole d'acc\`es, reprise sur panne,
        journalisation 
  \item {\bf Concurrence d'acc\`es}: multi-utilisateurs, transactions,
        partage
  \item {\bf Optimisation}~: index, hachage, gestion des tampons,
	algorithmes divers, ... 
\end{itemize}

$\rightarrow$ on va s'int\'eresser au(x) mod\`ele(s) logique(s).
\end{slide}

%%%%%%%%%%%%%%%%%%%%%%%%%%%%%%%%%%%%%%%%%%%%%%%%%%

\begin{slide}
\subsection{Caract\'eristiques d'un mod\`ele logique}

Mod\`ele de donn\'ees = fa\c{c}on de repr\'esenter les informations
du monde r\'eel. Un mod\`ele logique doit poss\'eder les
propri\'et\'es suivantes~:

\begin{itemize}
  \item Il est bas\'e sur des {\bf structures} pour repr\'esenter l'information
	(graphes, arbres, tableaux, listes, ensembles, ...)
  \item Il est associ\'e \`a des {\bf Langages} pour: 
	\begin{itemize}
          \item {\bf D\'efinir le mod\`ele de donn\'ees}~: langage de
                d\'efinition de donn\'ees (LDD)
          \item {\bf Interroger} et {\bf modifier} les donn\'ees~: 
		langages de
                requ\^etes et de manipulation de donn\'ees (LMD)
	\end{itemize}
  \item Il est asbtrait (ind\'ependant de l'implantation  physique)
  \item Il est possible de convertir automatiquement la
        repr\'esentation logique en repr\'esentation physique (des
	fichiers)
\end{itemize}
	
%%%%%%%%%%%%%%%%%%%%%%%%%%%%%%%%%%%%%%%%%%%%%%%%%%

\begin{slide}[1.5cm]
  \subsection{Architecture ANSI-SPARC}
  \begin{itemize}
    \item Sch\'ema Externe~: d\'efinition des vues adapt\'ees aux
          utilisateurs  
          
    \item Sch\'ema Conceptuel~: d\'efinition des besoins de donn\'ees
          ``r\'eelles'' ind\'ependamment du SGBD

    \item Sch\'ema Interne: 

    \begin{itemize}
      \item Niveau Logique~: description de l'organisation des
            donn\'ees en utilisant des mod\`eles de donn\'ees 
            $\Rightarrow$ sch\'ema logique

      \item Niveau Physique~: gestion de la m\'emoire secondaire (fichiers
            et enregistrements), index  
            $\Rightarrow$ sch\'ema physique
    \end{itemize}
  \end{itemize}
\vfill{}
\newpage
\centerline{\psfig{height=\textheight,figure=fi-archie-ansi.ps}}
\end{slide}

%%%%%%%%%%%%%%%%%%%%%%%%%%%%%%%%%%%%%%%%%%%%%%%%%%
%%%%%%%%%%
%%%%%%%%%%%%%%%%%%%%%%%%%%%%%%%%%%%%%%%%%%%%%%%%%%
%%%%%%%%%%
\begin{slide}[1.5cm]
\subsection{Ind\'ependance des Donn\'ees}
\begin{itemize}
  \item Ind\'ependance Physique: 
        
        On peut modifier l'implantation physique sans modifier le 
        niveau logique.\\
        
        Exemple: Changement/ajout d'un index ne modifie pas le sch\'ema
        logique dans la base de donn\'ees.
   
  \item Ind\'ependance Logique:
        
        On peut modifier le niveau logique sans toucher les
	applications existantes.\\
        
        Exemple: l'ajout d'une relation ne change pas le code 
	des applications.
        
 \end{itemize}
\end{slide}


%%%%%%%%%%%%%%%%%%%%%%%%%%%%%%%%%%%%%%%%%%%%%%%%%%
%%%%%%%%%%
%%%%%%%%%%%%%%%%%%%%%%%%%%%%%%%%%%%%%%%%%%%%%%%%%%
%%%%%%%%%%
\begin{slide}[1.5cm]
  \subsection{Historique des Mod\`eles de Donn\'ees}
  \begin{itemize}
    \item Ann\'ees 60-70~: mod\`eles hi\'erarchiques et r\'eseaux avec
          des langages navigationnels.
    \item Ann\'ees 80~: mod\`ele relationnel avec des langages d\'eclaratifs.
    \item Ann\'ees 90~: nouvelles applications $\rightarrow$ mod\`eles
          relationnels \'etendus, \`a valeurs structur\'ees et 
	  orient\'ees-objet.
  \end{itemize}
\vfill{}
\end{slide}

%%%%%%%%%%%%%%%%%%%%%%%%%%%%%%%%%%%%%%%%%%%%%%%%%%
%%%%%%%%%%
\begin{slide}[4cm]
\subsection{Exemple: Officiel de Spectacle}
\begin{itemize}
  \item Entit\'es: Cin\'ema, Salle, Film, Acteur, Directeur, Artiste
  \item Attributs, Clefs~: nom, titre, ville
  \item Associations~: s\'eance, joue, dirige
  \item H\'eritage~: est-un(e)
\end{itemize}
\end{slide}

%%%%%%%%%%%%%%%%%%%%%%%%%%%%%%%%%%%%%%%%%%%%%%%%%%
%%%%%%%%%%
\begin{slide}
\centerline{\psfig{width=\textwidth,figure=exempleER.ps}}
\end{slide}

%%%%%%%%%%%%%%%%%%%%%%%%%%%%%%%%%%%%%%%%%%%%%%%%%%
%%%%%%%%%%
\section{Le Mod\`ele Relationnel}

%%%%%%%%%%%%%%%%%%%%%%%%%%%%%%%%%%%%%%%%%%%%%%%%%%
%%%%%%%%%%
\begin{slide}
\subsection{Rappels sur le mod\`ele relationnel}
 
\begin{enumerate}
  \item Repr\'esentation des entit\'es par des {\bf tables}. Chaque ligne
        est une occurence de l'entit\'e (un n-uplet), chaque colonne un
        attribut.
  \item Les attributs sont d\'efinis sur des types de base (entiers,
        cha\^{\i}nes de caract\`eres)
  \item Le sch\'ema d'une relation est compos\'e d'un nom et d'une
	liste de noms d'attributs avec leurs domaines.
  \item Interrogation par SQL

\end{enumerate}
 
\end{slide}
 

%%%%%%%%%%%%%%%%%%%%%%%%%%%%%%%%%%%%%%%%%%%%%%%%%%
%%%%%%%%%%%%%%%%%%
\begin{slide}
\subsection{Relation~: Exemples}

\begin{center}
\begin{tabular}{||c|c|c||}\hline
\multicolumn{3}{||c||}{CIN\`EMA}\\\hline Cin\'ema &Adresse & Ville\\ \hline
"Rex"  & ``7, Bd Bonne Nouvelle'' & ``PARIS'' \\ \hline
"Cinefil"  & ``8, rue Dunod'' & ``CRETEIL'' \\ \hline
"Kino" &  ``123, Bd Raspail'' & ``PARIS''\\ \hline
...  &...  &...  \\ \hline
\end{tabular}
\end{center}
 
\vspace*{1cm}

\begin{tabular}{||c|c|c||}\hline
\multicolumn{3}{||c||}{SALLE}\\\hline 
Cin\'ema & No salle & Capacit\'e \\ \hline
``Rex'' & 1 & 200 \\\hline 
``Rex'' & 3 & 120 \\\hline 
``Kino'' & 1 & 300\\\hline 
... &...  &...  \\\hline
\end{tabular}
\hfill
\begin{tabular}{||c|c|c|c||}\hline
\multicolumn{4}{|||c|||}
{SEANCE}\\\hline
Cinema & Salle & Film & Horaire\\\hline
``Rex'' & 1 & ``Cyrano'' & 16-18 \\\hline 
``Kino'' & 3 & ``Les oiseaux'' & 14-16 \\\hline
...  &... \\\hline
\end{tabular}
\end{slide}


%%%%%%%%%%%%%%%%%%%%%%%%%%%%%%%%%%%%%%%%%%%%%%%%%%
%%%%%%%%%%%%%%%%%%
\begin{slide}[1cm]
\subsection{Manipulation et Interrogation des Donn\'ees}
\begin{itemize}
\item Langage de D\'efinition des Donn\'ees (LDD)~: cr\'eation et
      modification de sch\'emas
\item Langage de Manipulation des Donn\'ees (LMD)~: cr\'eation, modification
      et suppressions de n-uplets
\item Langage d'Interrogation des Donn\'ees (LR)~: cr\'eation de nouvelles
      relations (intentionnelles) \`a partir des relations existantes
\end{itemize}
\vfill
\end{slide}
%%%%%%%%%%%%%%%%%%%%%%%%%%%%%%%%%%%%%%%%%%%%%%%%%%
%%%%%%%%%%%%%%%%%%
%%%%%%%%%%%%%%%%%%%%%%%%%%%%%%%%%%%%%%%%%%%%%%%%%%
%%%%%%%%%%%%%%%%%%
%\begin{slide}[1cm]
%\subsection{Langages de Requ\^etes: Alg\`ebre et Calcul} 
%
%On peut distinguer entre deux types de langages formels~:
%\begin{itemize}
%\item Calcul Relationnel~: une requ\^ete correspond \`a un pr\'edicat
%      dans la logique du premier ordre
%\item Alg\`ebre Relationnelle~: \ 
%	\begin{itemize}
%	\item op\'erations de base~: s\'election ($\sigma$), projection
%              ($\pi$), produit cart\'esien ($\times$), renommage~($\rho$)
%	\item op\'erations ensemblistes~: union, diff\'erence, intersection
%	\item op\'erations d\'eriv\'es~: jointure, division, \ldots{}
%	\end{itemize}
%\item Le langage SQL est issu du calcul relationnel
%\end{itemize}
%\end{slide}
%
%%%%%%%%%%%%%%%%%%%%%%%%%%%%%%%%%%%%%%%%%%%%%%%%%%
%%%%%%%%%%%%%%%%%%
%
%\begin{slide}[7cm]
%\subsection{Exemples de Requ\^etes I}
%
%\centerline{\Underline{"Tous les films qui passent au REX" }}
%
%\begin{itemize}
%\item Alg\`ebre :
% 
% $\pi_{Film} (\sigma_{Cin\'ema=''Rex''}(SEANCE))$
%
%\item Calcul Relationnel:
%
% $\{f | \exists x \exists y \exists z: Seance(x,y,f,z) \wedge x = ``Rex''\}$
%
%\item SQL: \
%\begin{verbatim}
%      select film
%        from SEANCE
%       where Nom-cinema = ``Rex''
%\end{verbatim}
%\end{itemize}
%\end{slide}

%%%%%%%%%%%%%%%%%%%%%%%%%%%%%%%%%%%%%%%%%%%%%%%%%%

\begin{slide}[7cm]
\subsection{Exemples de Requ\^etes}

\centerline{\Underline{"Quels films passent \`a Cr\'eteil ?"}}

\begin{itemize}

\item Il faut faire une {\bf jointure} entre SEANCE et CINEMA~: \
\begin{verbatim}
      select film
        from SEANCE, CINEMA
        where CINEMA.nom = SEANCE.nom-cinema
        and	CINEMA.ville = ``CRETEIL''
\end{verbatim}
\end{itemize}
\end{slide}

%%%%%%%%%%%%%%%%%%%%%%%%%%%%%%%%%%%%%%%%%%%%%%%%%%
%%%%%%%%%%%%%%%%%%

\begin{slide}
\subsection{Avantages du Mod\`ele Relationnel}
\begin{itemize}
\item Simplicit\'e (une seule structure)
\item D\'eclarativit\'e (compl\`ete ind\'ependance logique/physique)
\item Facilit\'e d'emploi
\item SQL
\item C'est le standard actuel...
\end{itemize}
\vfill{}
\end{slide}

%%%%%%%%%%%%%%%%%%%%%%%%%%%%%%%%%%%%%%%%%%%%%%%%%%

\begin{slide}[1.5cm]
\subsection{Inconv\'enients du Mod\`ele Relationnel}
\begin{itemize}
\item {\bf Pauvret\'e du mod\`ele}~: impossibilit\'e de repr\'esenter
	des donn\'ees complexes, pas de notion d'{\bf identifiant} 
	d'une entit\'e

\item {\bf Pauvret\'e du langage}~:
  \begin{enumerate}
	\item  On ne peut pas exprimer toutes 
	les requ\^etes $\rightarrow$ int\'egration avec des langages de
	programmation (C+SQL)
	\item S\'eparation donn\'ees/traitements~: pas d'association
	entre les aspects dynamiques et statiques des donn\'ees
  \end{enumerate}

\item {\bf Probl\`emes de performances} (op\'erations ensemblistes)

\end{itemize}
\vfill{}
\end{slide}

%%%%%%%%%%%%%%%%%%%%%%%%%%%%%%%%%%%%%%%%%%%%%%%%%%
%%%%%%%%%%%%%%%%%%

\begin{slide}[1.5cm]
\subsection{Pauvret\'e de la Repr\'esentation}

Une seule structure~: la table.

\begin{itemize}
\item Perte de s\'emantique~: on ne distingue plus les
	entit\'es et les associations (Ex~: S\'eance)

\item Repr\'esentation d'objets complexes par plusieurs tables (Ex~:
	Cin\'ema, Salles, S\'eance). 
	
\item Impossible de distinguer des groupes d'attributs dans une
	relation (Ex~: attribut ``Adresse'' dans Cin\'ema).
	
\item Pas de notion ``d'ensemble de valeurs''. Exemple~: chaque
	horaire est un n-uplet alors que son importance s\'emantique
	est faible.

\item Pas de notion d'identit\'e. Elle doit \^etre g\'er\'ee
	par le programmeur.

\item Typage faible~: on voudrait d\'efinir un type ``P\'eriode
	de temps'' par exemple.
\end{itemize}

\end{slide}

%%%%%%%%%%%%%%%%%%%%%%%%%%%%%%%%%%%%%%%%%%%%%%%%%%
%%%%%%%%%%%%%%%%%%

\begin{slide}[1.5cm]
\subsection{Pauvret\'e du Langage~ LP}

\begin{enumerate}
   \item Pas de r\'ecursivit\'e. \\
	Exemple~: {\it Quelles sont toutes
		les suites d'un film~?}
   \item Langage non proc\'edural\\
        Exemple~: repr\'esentation d'entit\'es spatiales~:
\begin{enumerate}
\item Point (id, x, y)
\item Ligne (id\_ligne, no\_point, id\_point)
\item Zone (id\_zone, no\_ligne, id\_ligne)
\end{enumerate}
 
On voudrait associer la fonction {\it Surface} \`a Zone,
 {\it Longueur} \`a Ligne, etc.\\
\end{enumerate}

$\rightarrow$ On combine SQL et un langage proc\'edural (le C par exemple)

\end{slide}

%%%%%%%%%%%%%%%%%%%%%%%%%%%%%%%%%%%%%%%%%%%%%%%%%%

\begin{slide}[1.5cm]
\subsection{Probl\`emes de performances}

\centerline{\Underline{"Quels sont les points formant le contour de la zone 2?"}}

Il faut faire une {\bf jointure} entre Zone, Ligne et Point~: \\
\begin{verbatim}
      select x,y
        from Zone, Ligne, Point
        where Zone.id-zone=2
        and   Zone.id_ligne=Ligne.id-ligne
        and   Ligne.id-point=Point.id-point
\end{verbatim}

Si un zone est d\'ecrite par 50 lignes, et une ligne par 100 points,
il faudra 5~000 jointures pour reconstituer le contour~!

\end{slide}

%%%%%%%%%%%%%%%%%%%%%%%%%%%%%%%%%%%%%%%%%%%%%%%%%%
%%%%%%%%%%%%%%%%%%

\section{Du Relationnel vers les SGBD Objets}

%%%%%%%%%%%%%%%%%%%%%%%%%%%%%%%%%%%%%%%%%%%%%%%%%%
%%%%%%%%%%%%%%%%%%

% Slide 20

\begin{slide}
\subsection{Du Relationnel vers les SGBD Objets}

Comment peut-on adapter les SGBD aux nouveaux besoins~?  
\begin{itemize}
  \item Plus de structures~: ensembles, listes, tableaux ...

  \item Plus de manipulations complexes

\end{itemize}

Quelques propositions~:

\begin{itemize}
	\item Mod\`eles \`a valeurs structur\'ees
	\item Mod\'eles relationnels \'etendus
	\item Mod\`eles orient\'es-objet
\end{itemize}
\hfill
\end{slide}


%%%%%%%%%%%%%%%%%%%%%%%%%%%%%%%%%%%%%%%%%%%%%%%%%%
%%%%%%%%%%

\begin{slide}
\subsection{Valeurs Structur\'ees~: relations N1NF}

Les constructeurs d'ensemble et de n-uplet sont appliqu\'es {\em
alternativement}. Ainsi, il n'est pas possible de repr\'esenter des
ensembles d'ensembles ou des n-uplets de n-uplets.

\begin{center}
  \begin{tabular}{|c|c|c|}
    \multicolumn{3}{c}{\bf FILMS}    \\\hline
    {\bf titre} & {\bf acteurs}      & {\bf cin\'emas} \\\hline
    Arizona Dream & Johnny Depp      & Forum Horizon \\
                  & Faye Dunaway     & Luxembourg \\
                  & Jerry Lewis      &            \\\hline
    Fatale        & Juliette Binoche & Montparnasse \\
                  & Jeremy Irons     & \\\hline
    \ldots{}      & \ldots{}         & \ldots{} \\\hline
  \end{tabular}
\end{center}

Probl\`eme~: pas d'identit\'e pour une valeur, donc pas possible
		de la {\bf partager}.
\end{slide}

%%%%%%%%%%%%%%%%%%%%%%%%%%%%%%%%%%%%%%%%%%%%%%%%%%

%\begin{slide}
%\subsection{\'Evolution du Relationnel \`a l'Objet}

%\centerline{\psfig{width=\textwidth,figure=fi-reet.ps}}

%\end{slide}

%%%%%%%%%%%%%%%%%%%%%%%%%%%%%%%%%%%%%%%%%%%%%%%%%%
%%%%%%%%%%

\begin{slide}
\subsection{Relationnel \'etendu: les Types Abstraits de Donn\'ees}

\Underline{Type Abstrait = Structure de Donn\'ees + Ensemble
d'Op\'erations}

Les donn\'ees sont {\bf encapsul\'ees} et on y acc\`ede
uniquement \`au travers des op\'erations d\'efinies
dans le type

\begin{verbatim}
     type Point     = [      type Rectangle = [
          x:float,                xmin:float, ymin:float,
          y:float]                xmax:float, ymax:float]

     function contains(Rectangle, Point): Booleen
     create table RECTANGLES(nom: string, geometry: Rectangle)

     select R.nom
       from R in RECTANGLES
      where contains(R.geometry,"10.4,23.5")
\end{verbatim}

\end{slide}

%%%%%%%%%%%%%%%%%%%%%%%%%%%%%%%%%%%%%%%%%%%%%%%%%%
%%%%%%%%%%

\section{Les Mod\`eles Orient\'es-Objet}

\begin{slide}
  \subsection{SGBD Orient\'es-objet}
 
  Principles caract\'eristiques~:
 
  \begin{enumerate}
    \item Fusion des technologies SGBD et langages de programmation
    \item Structuration plus riche qu'en relationnel
    \item Association des donn\'ees et des traitements
  \end{enumerate}
 
  Bibliographie~:
 
  \begin{enumerate}
      \item Delobel, L\'ecluse, Richard, 
	{\bf Bases de donn\'ees~: des syst\`emes
                relationnels aux syst\`emes \`a objets}
      \item Adiba, Collet, {\bf Le SGBD O$_{2}$}
      \item Catell (ed.), {\bf The Object Database Standard}
   \end{enumerate}
\end{slide}

%%%%%%%%%%%%%%%%%%%%%%%%%%%%%%%%%%%%%%%%%%%%%%%%%%
%%%%%%%%%%

\begin{slide}[1.5cm]
\subsection{Objet: Identificateur et Valeur Structur\'ee}

Un objet, c'est une paire~:

\begin{enumerate}
  \item identificateur ({\it oid})
  \item valeur
\end{enumerate}

Caract\'eristiques~:

\begin{enumerate}
  \item L'objet garde le m\^eme {\it oid} toute sa vie
  \item L'{\it oid} est ind\'ependant de la valeur de l'objet
  \item La valeur est d\'ecrite par un {\it type} qui peut \^etre
	complexe
\end{enumerate}
\end{slide}

\begin{slide}

%%%%%%%%%%%%%%%%%%%%%%%%%%%%%%%%%%%%%%%%%%%%%%%%%%
%%%%%%%%%%

\subsection{Exemples d'Objets}

\begin{verbatim}
(i0, [Titre: 'Rambo', Acteurs: {i3,i4}, Suites: {i1,i2}])
(i1, [Titre: 'Rambo II', Acteurs: {i3,i9}, Suites: {i2}])
\end{verbatim}

Autre notation~:
\begin{verbatim}
          (i3, tuple[Nom:`Stallone',
                      adresse : tuple [No : 40,
                                       Rue : ``Rambo Street'',
                                       Ville : berverly-hills]
                      Films : {i0, i1, i2}
                     ]
          )
\end{verbatim}

Les valeurs des objets sont construites \`a partir de valeurs
atomiques et d'identificateurs d'autres objets en utilisant
les constructeurs ensemble et tuple (n-uplet) dans n'importe
quel ordre.
\end{slide}

%%%%%%%%%%%%%%%%%%%%%%%%%%%%%%%%%%%%%%%%%%%%%%%%%%

\begin{slide}
\subsection{Composition d'objets}

L'identit\'e permet~:

\begin{itemize}
  \item La {\bf composition} d'un objet avec d'autres objets
  \item Le {\bf partage} d'un m\^eme objet entre plusieurs objets compos\'es
\end{itemize}

Dans les exemples pr\'ec\'edents~:

\begin{itemize}
  \item L'objet film 'Rambo' est compos\'e de l'objet acteur 'Stallone'
  \item L'objet acteur 'Stallone' est composant de plusieurs films
\end{itemize}

$\rightarrow$ La composition permet de 
	 repr\'esenter avec le mod\`ele OO les associations
 	du mod\`ele E/R.
\end{slide}

%%%%%%%%%%%%%%%%%%%%%%%%%%%%%%%%%%%%%%%%%%%%%%%%%%

\begin{slide}
\subsection{Identit\'e et \'Egalit\'e}
\begin{itemize}
\item Deux objets de m\^eme identificateur sont {\bf identiques}\\
	
   	(i,1) et (i,3) repr\'esentent le m\^eme objet\\

\item Deux objets sont {\bf \'egaux} s'ils ont la m\^eme valeur~: \\

	$(i,5)$ est \'egal \`a $(j,5)$\\ 

	$(i,[A: k, B: 5])$ est \'egal \`a $(j,[A:k,B:5])$.
\end{itemize}
\vfill
\end{slide}

%%%%%%%%%%%%%%%%%%%%%%%%%%%%%%%%%%%%%%%%%%%%%%%%%%

\begin{slide}[1.5cm]
\subsection{\'Egalit\'e en Profondeur}

\Underline{D\'efinition:}
\begin{itemize}
\item Deux objets \'egaux sont \'egaux en profondeur
\item Deux objets qui ont les m\^emes attributs et dont les valeurs
      sont
  \begin{itemize}
    \item soit \'egales
    \item soit des identificateurs d'objets \'egaux en
          profondeur
  \end{itemize}
sont \'egaux en profondeur.
\end{itemize}
\end{slide}

%%%%%%%%%%%%%%%%%%%%%%%%%%%%%%%%%%%%%%%%%%%%%%%%%%

\begin{slide}
\subsection{Egalit\'e, \'egalit\'e en profondeur}
 
\begin{verbatim}
      o1: (i1,[A:i3,B:b])   
      o2: (i2,[A:i4,B;b])
      o3: (i3,5)            
      o4: (i4,5)
\end{verbatim}

\begin{enumerate}
  \item L'objet o1 n'est pas \'egal \`a o2
  \item L'objet o3 est \'egal \`a o4
  \item o1 est \'egal en profondeur \`a o2.
\end{enumerate}
\end{slide}

%%%%%%%%%%%%%%%%%%%%%%%%%%%%%%%%%%%%%%%%%%%%%%%%%%
%%%%%%%%%%

\begin{slide}
\subsection{Objets et Classes}

Un objet est d\'efini par~:

\begin{itemize}
  \item Son {\it oid}
  \item Sa valeur 
  \item Sa structure
  \item Son comportement (= ensemble des fonctions qu'on peut
	appliquer \`a sa valeur)
\end{itemize}

3 et 4 constituent la {\bf classe} de l'objet. Une classe est donc 
la description de tous les objets ayant m\^eme structure et m\^eme 
comportement.

\noindent
{\bf Attention} : l'ensemble des objets d'une classe n'est pas
forc\'ement stock\'e dans la base.\\

\end{slide}

%%%%%%%%%%%%%%%%%%%%%%%%%%%%%%%%%%%%%%%%%%%%%%%%%%
%%%%%%%%%%

\begin{slide}
\subsection{Classe : le type}

Un type (structure) est construit \`a partir~:

\begin{itemize}
\item D'un ensemble de types atomiques (entier, r\'eel, caract\`ere,
cha\^\i ne,...)
\item D'un ensemble $C$ de noms de classes
\item De constructeurs (n-uplet, ensemble, liste, tableau, ...)
\end{itemize}

\begin{enumerate}
\item Un type atomique est un type.
\item Un \'el\'ement de $C$ est un type.
\item Si $t1,...,tn$ sont des types et $a1,...,an$ sont des attributs,
      alors $[a1:t1,...,an:tn]$ est un type.
\item Si $t$ est un type, alors $\{t\}$ est un type.
\end{enumerate}

Cette d\'efinition de type est celle des valeurs structur\'ees \`a
laquelle on rajoute des noms de classes comme types.
\end{slide}

%%%%%%%%%%%%%%%%%%%%%%%%%%%%%%%%%%%%%%%%%%%%%%%%%%
%%%%%%%%%%

\begin{slide}[1cm]
\subsection{Classe : le Comportement}

Les m\'ethodes d'un type d'objet sont l'ensemble des fonctions
(proc\'edures) applicables aux objets ayant ce type.\\

Le comportement d'un objet est d\'efini par un
ensemble de signatures de m\'ethodes.\\
 
La signature d'une m\'ethode et compos\'ee
\begin{enumerate}
\item du nom de la m\'ethode
\item du type des arguments
\item du type du r\'esultat
\end{enumerate}

\Underline{Exemple:} type PERSONNE avec m\'ethode {\it \^age~}: INT qui
calcule l'\^age \`a partir de la date de naissance.

\end{slide}

%%%%%%%%%%%%%%%%%%%%%%%%%%%%%%%%%%%%%%%%%%%%%%%%%%

\begin{slide}
\subsection{Sch\'ema d'une Base de Donn\'ees}

Un sch\'ema est un ensemble de classes et types d'objets. Exemple~:

\begin{verbatim}
         class PERSONNE: 
            type tuple (�poux: PERSONNE,
                        ss_nom: tuple (ss:INT,nom:STRING)
                        enfants: set(PERSONNE))
            method �ge: INT

         class ENTREPRISE:
            type tuple   (nom: STRING,
                     employ�s: set(PERSONNE))
\end{verbatim}
\end{slide}

%%%%%%%%%%%%%%%%%%%%%%%%%%%%%%%%%%%%%%%%%%%%%%%%%%

\begin{slide}[3cm]
\subsection{Concepts OO : Message}

Seule mani\`ere d'ex\'ecuter un traitement: envoi d'un message \`a un 
objet dont le type contient la m\'ethode \`a ex\'ecuter

\Underline{message:} objet$\rightarrow$m\'ethode(a1,...,an)

\Underline{Exemple:} fran\c{c}ois est un objet de type
PERSONNE. L'envoi du message fran\c{c}ois$\rightarrow$ \^age \`a
fran\c{c}ois donne son \^age.  \vfill
\end{slide}

%%%%%%%%%%%%%%%%%%%%%%%%%%%%%%%%%%%%%%%%%%%%%%%%%%

\begin{slide}
\subsection{Concepts OO : Encapsulation}

Un objet est caract\'eris\'e par sa structure et par son
interface. L'interface est un ensemble de m\'ethodes visible par
l'ext\'erieur. Un objet ne peut \^etre manipul\'e que via son
interface $\Rightarrow$ il s'agit du principe d'encapsulation. \`A
l'ext\'erieur,

\begin{itemize}
\item On voit la signature des m\'ethodes 
\item On ne voit pas la structure des objets  ni le
	corps des m\'ethodes.
\end{itemize}

\Underline{Programmation Modulaire~:}

\begin{itemize}
\item S\'eparation entre l'implantation et la sp\'ecification
\item Regroupement dans un m\^eme module des donn\'ees et de leurs
      traitement
\item Possibilit\'e d'avoir plusieures implantations pour une m\^eme
      sp\'ecification 
\end{itemize}
\end{slide}

%%%%%%%%%%%%%%%%%%%%%%%%%%%%%%%%%%%%%%%%%%%%%%%%%%

\begin{slide}[1.5cm]
\subsection{Traitements sur les Donn\'ees}
\begin{itemize}
\item Relationnel: \
\begin{itemize}
\item appel SQL \`a la base dans le programme d'applications
\item donn\'ees dans la base et programmes dans une biblioth\`eque
\item conception et sp\'ecification des programmes sont
      ind\'ependantes des donn\'ees 
\item toute donn\'ee est accessible \`a tout programme
\end{itemize}

\item SGBD Orient\'e-Objet Int\'egr\'e: \
\begin{itemize}
\item m\'ethodes
\item donn\'ees et programmes dans la base
\item donn\'ees seulement accessibles par des m\'ethodes de la m\^eme
      classe (public, private, protected)
\end{itemize}
\end{itemize}
\end{slide}

%%%%%%%%%%%%%%%%%%%%%%%%%%%%%%%%%%%%%%%%%%%%%%%%%%

\begin{slide}
\subsection{H\'eritage: Classes - Sous-Classes}

Une classe $C$ est une \Underline{sous-classe} d'une classe $C'$
(h\'erite de $C'$), si
\begin{itemize}
\item chaque instance de (objet dans) $C$ est une instance (un
      objet dans) $C'$.
\item toutes les m\'ethodes de $C$ peuvent \^etre appliqu\'ees aux
      instances de $C'$.
\end{itemize}

\Underline{Exemple:}
\begin{verbatim}
        class PERSONNE: type [nom:STRING, date: DATE, sexe: BOOL]
              method �ge(DATE): INT

        class EMPLOY�: type [nom:STRING, date: DATE, sexe: BOOL
              method dipl�mes: {STRING}
                     �ge(DATE): INT
                     salaire:INT]
\end{verbatim}

\end{slide}

%%%%%%%%%%%%%%%%%%%%%%%%%%%%%%%%%%%%%%%%%%%%%%%%%%

\begin{slide}[1.5cm]
\subsection{H\'eritage : polymorphisme et surcharge}

\Underline{Exemple :} classes g\'eom\'etriques

\begin{verbatim}
  Class Zone      Methodes : Intersection, Surface, ...
    Class Rectangle inherit Zone
    Class Cercle    inherit Zone
    Class Triangle  inherit Zone
\end{verbatim}
            
{\bf Polymorphisme}~: Le m\^eme code peut s'appliquer \`a des objets de
	classes diff\'erentes

{\bf Surcharge}~: On peut red\'efinir les m\'ethodes pour chaque sous-classe.\\

\vfill{}
\end{slide}

%%%%%%%%%%%%%%%%%%%%%%%%%%%%%%%%%%%%%%%%%%%%%%%%%%

\begin{slide}
\subsection{H\'eritage~: r\'eutilisation}

\centerline\Underline{On veut calculer la surface de diff\'erents
objets g\'eom\'etriques:}

\begin{verbatim}
   sans objets                          avec objets
   -----------                          -----------
   for all g in geo_set do              for all g in geo_set
     case of type(g):                     g->surface
        cercle: surface_cercle(g)
      triangle: surface_triangle(g)
     rectangle: surface_rectangle(g)
\end{verbatim}

\end{slide}
%%%%%%%%%%%%%%%%%%%%%%%%%%%%%%%%%%%%%%%%%%%%%%%%%%

\begin{slide}
\subsection{Typage dynamique}

Le mod\`ele OO implique un certain degr\'e de typage dynamique~:

\noindent
{\bf R\'esolution tardive} : le syst\`eme d\'etermine 
    {\bf au moment de l'ex\'ecution} la classe de l'objet et donc 
    la m\'ethode \`a appliquer

\noindent
{\bf G\'en\'ericit\'e}
\begin{itemize}
  \item Le programmeur n'a pas besoin de conna\^\i tre le type de l'objet
        dont il calcule la surface.
       
  \item L'adjonction d'un nouveau type d'objet set fait sans
        modification du programme appelant.

\end{itemize}
\end{slide}

%%%%%%%%%%%%%%%%%%%%%%%%%%%%%%%%%%%%%%%%%%%%%%%%%%

\begin{slide}[1.5cm]
\subsection{H\'eritage Multiple}
\begin{verbatim}
       class PERSONNE = type [nom: STRING, ss: INT, naissance: INT]
             method �ge: INT

       class EMPLOY� inherits PERSONNE = type [salaire:INT]
             method augmente_salaire: INT

       class �TUDIANT inherits PERSONNE = type [cours:{COURS}]
             method notes: {INT}

       class ATER inherits EMPLOY�, �TUDIANT = type [heures:{INT}]
             method temps_recherche: INT
\end{verbatim}

 
ATER h\'erite des attributs et m\'ethodes de ses 2 superclasses 
EMPLOY\'E et \'ETUDIANT.
\end{slide}

%%%%%%%%%%%%%%%%%%%%%%%%%%%%%%%%%%%%%%%%%%%%%%%%%%

\begin{slide}[1.5cm]
\subsection{Conflit d'H\'eritage }

\begin{verbatim}
       class PROFESSEUR inherits PERSONNE = 
             type [donne:{COURS}, horaires:HEURES]

       class �TUDIANT inherits PERSONNE = 
             type [suit:{COURS}, horaires:HEURES]

       class ASSISTANT inherits PROFESSEUR, �TUDIANT 
             rename horaires in PROFESSEUR:
                    horaires_enseignant = 
             type [temps_recherche: HEURES] 
\end{verbatim}


ASSISTANT h\'erite deux attributs de m\^eme nom (horaires) $\Rightarrow$
il faut renommer l'un d'eux (ici celui h\'erit\'e de la classe
PROFESSEUR)
\end{slide}

%%%%%%%%%%%%%%%%%%%%%%%%%%%%%%%%%%%%%%%%%%%%%%%%%%

\begin{slide}[1.5cm]
\subsection{Utilisation de l'H\'eritage}
\begin{itemize}
\item {\bf R\'eutilisation de code}~: les m\'ethodes d'une classe sont
      r\'eutilis\'ees dans ses sous-classes: la m\'ethode \^age a le
      m\^eme code pour un employ\'e et pour une personne (diff\'erence
      entre l'ann\'ee actuelle et l'ann\'ee de la naissance)

\item {\bf Surcharge et liens dynamiques}~: la m\'ethode {\it \^age} est
      red\'efinie pour les employ\'es. Une m\'ethode surcharg\'es a le
      m\^eme nom que la m\'ethode dans la super-classe mais un code
      diff\'erent.  Le ``bon'' code est trouv\'e \`a l'ex\'ecution par
      des liens dynamiques (r\'esolution tardive).

\item {\bf Modularit\'e}~: les modifications dans les m\'ethodes d'une classe
      change uniquement le comportement de ses sous-classes.
\end{itemize}
La surcharge et la r\'eutilisation permettent un certain degr\'e de 
polymorphisme. 
\end{slide}

%\sectioncontents
\section{La persistance dans les SGBDOO}

%%%%%%%%%%%%%%%%%%%%%%%%%%%%%%%%%%%%%%%%%%%%%%%%%%

\begin{slide}[1.5cm]
\subsection{Probl\'ematique de la persistance}

Un SGBDOO doit permettre une int\'egration compl\`ete des
outils de programmation (types, langages) et des outils
'Base de donn\'ees' (stockage, extraction)\\

Un objet doit donc pouvoir \^etre~:

\begin{itemize}
\item {\bf Temporaire} : c'est une 'variable' dans un programme. Il
	n'existe que pendant l'ex\'ecution.
\item {\bf Persistant} : il existe avant, pendant et apr\`es
	l'ex\'ecution d'un programme.
\end{itemize}

D'o\`u le probl\`eme~: {\bf comment d\'eclare-t-on dans un programme 
	les donn\'ees persistantes}

\end{slide}

%%%%%%%%%%%%%%%%%%%%%%%%%%%%%%%%%%%%%%%%%%%%%%%%%%

\begin{slide}[1.5cm]
\subsection{Quelques qualit\'es du mod\`ele de persistance}

\begin{itemize}
\item {\bf Orthogonalit\'e au type} : toute classe doit pouvoir contenir
	des objets temporaires ou persistants.\\
	Contre-exemple : les SGBDR.

\item {\bf Persistance \`a la cr\'eation ou en cours de programme} : 
	une donn\'ee temporaire doit pouvoir devenir persistante \`a tout
	moment.

\item {\bf Compatibilit\'e du code source} : le m\^eme programme doit
	pouvoir manipuler des donn\'ees temporaires ou persistantes
	de mani\`ere transparente.
\end{itemize}

\end{slide}


%%%%%%%%%%%%%%%%%%%%%%%%%%%%%%%%%%%%%%%%%%%%%%%%%%

\begin{slide}[1.5cm]
\subsection{Persistance par attachement}

Ne sont persistants :

\begin{itemize}
\item Que les points d'entr\'ee ou {\bf Racines de persistance}.

\item Que les objets accessibles par composition \`a partir
	d'un point d'entr\'ee.
\end{itemize}

Les points d'entr\'ee sont d\'efinis par des {\bf valeurs nomm\'ees}.
\begin{verbatim}
        NAME films     : set(Film)
        NAME artistes  : set(Artiste)
        NAME Belmondo  : Acteur
\end{verbatim}

Le m\^eme objet (Belmondo) peut \^etre accessible \`a partir de plusieurs
points d'entr\'ee.

\end{slide}

%%%%%%%%%%%%%%%%%%%%%%%%%%%%%%%%%%%%%%%%%%%%%%%%%%

\section{Le langage de requ\^ete}

\begin{slide}[1.5cm]
\subsection{Extraction : le langage de requ\^ete OQL}

SQL n'est plus suffisant (car li\'e aux relations en premi\`ere
forme normale). Les nouveaux besoins sont~:

\begin{itemize}
\item Exprimer des crit\`eres et des s\'elections sur des valeurs complexes.

\item Parcourir un graphe de composition

\item Ex\'ecuter des fonctions ou des m\'ethodes

\end{itemize}

\end{slide}

%%%%%%%%%%%%%%%%%%%%%%%%%%%%%%%%%%%%%%%%%%%%%%%%%%

\begin{slide}[1.5cm]
\subsection{Exemple de sch\'ema I}

\begin{verbatim}
   class Cinema 
         tuple (nom : string,
              adresse : tuple (No : integer,
                               Rue : string,
                               ville : Ville)
              salles : set(Salle))

   class Salle
         tuple (No : integer,
                capacite: integer,
                seances : list(tuple(heure:Heure,
                                     film : Film))
                )
\end{verbatim}
\end{slide}

%%%%%%%%%%%%%%%%%%%%%%%%%%%%%%%%%%%%%%%%%%%%%%%%%%

\begin{slide}[1.5cm]
\subsection{Exemple de sch\'ema II}

\begin{verbatim}
    class Film
          tuple (titre : string,
                 acteurs: set(Acteur),
                 metteur_en_scene : Artiste)

    class Artiste
          tuple (nom:String)

    class Acteur inherit Artiste
          tuple (roles : set(tuple(film:Film,
                                  nom: string))
                )
\end{verbatim}

\end{slide}

%%%%%%%%%%%%%%%%%%%%%%%%%%%%%%%%%%%%%%%%%%%%%%%%%%

\begin{slide}[1.5cm]
\subsection{Exemple de requ\^etes I}

{\it Tous les films de Jean-Luc Godard}

\begin{verbatim}
          SELECT film.titre
          FROM   film in Films
          WHERE  film.metteur_en_scene.nom = ``Godard''
\end{verbatim}

\end{slide}

%%%%%%%%%%%%%%%%%%%%%%%%%%%%%%%%%%%%%%%%%%%%%%%%%%

\begin{slide}[1.5cm]
\subsection{Exemple de requ\^etes II}

{\it Quels acteurs ont jou\'e ``Hamlet''}

\begin{verbatim}
    SELECT acteur.nom
    FROM   acteur in Acteurs,
              role in acteur.roles
    WHERE  role.nom = ``Hamlet''
\end{verbatim}
\end{slide}

%%%%%%%%%%%%%%%%%%%%%%%%%%%%%%%%%%%%%%%%%%%%%%%%%%

\begin{slide}[1.5cm]
\subsection{Exemple de requ\^etes III}

{\it O\`u peut-on voir ``A bout de souffle''}

\begin{verbatim}
   SELECT acteur.nom
   FROM cinemas in Cinemas,
          salle in cinema.salle
            seance in salle.seances
   WHERE   seance.film.titre = ``A bout de souffle''
\end{verbatim}
\end{slide}

%%%%%%%%%%%%%%%%%%%%%%%%%%%%%%%%%%%%%%%%%%%%%%%%%%

\begin{slide}[1.5cm]
\subsection{Conclusion : situation actuelle des SGBDOO}

\begin{itemize}
  \item Une technologie r\'ecente, bas\'ee sur des concepts
	issus en grande partie du monde de la recherche
  \item Des SGBD complets existent (ObjectStore, Versant,
	O$_{2}$, Matisse)
  \item Les applications restent pour l'instant tr\`es sp\'ecifiques
	(Syst\`emes d'Information G\'eographiques, Multim\'edia,
	 Documents structur\'es, CAO)
\end{itemize}

Le march\'e d\'ecidera !

\end{slide}

\end{document}
